\subsection*{Why a bounded scalar primitive error cannot supply the finite floor}

The weaker target in Corollary~\ref{cor:v136-complement-floor} makes it
necessary to revisit the scalar primitive route: its earlier failure to
give an $o(1)$ error does not itself exclude an $O(1)$ error. The following
argument rules out that weaker scalar target for the \emph{actual}
arithmetic symbol. It does not rule out the signed concentration operator.
Keep $a=\log\lambda$, $T=2a$, the exact $\beta_a$ of
\eqref{eq:v130-beta}, and $\Delta_a=\Delta(\beta_a)$ from
\eqref{eq:v134-drawdown}.

\begin{lemma}[A positive probe with a negative truncated Fourier mass]
\label{lem:v137-probe}
Fix $\delta=1/100$ and an even nonnegative
$\kappa\in C_c^\infty(-\delta,\delta)$ with integral one. Put
\[
 w=\mathbf1_{[\pi,3\pi]}*\kappa,\qquad
 H(t)=\int_{\mathbb R}w(s)e^{-its}\,ds,\qquad
 K(u)=\int_{-1}^1H(t)e^{iut}\,dt.
\]
Then $w\ge0$, $\int w=2\pi$, $\norm{w'}_1=2$, and
\begin{equation}\label{eq:v137-probe-properties}
 H(1)=H(-1)=0,\qquad
 K(0)<-\frac1{3\pi},\qquad
 |K(u)|\le\frac{C_w}{1+u^2}.
\end{equation}
Here $K$ is real on the real line; no symmetry of $w$ about zero is
required.
\end{lemma}
\begin{proof}
The two translates in $w'(s)=\kappa(s-\pi)-\kappa(s-3\pi)$ have
disjoint supports, giving the stated variation. With the nonunitary
transform used only in this lemma,
\[
 H(t)=2e^{-2\pi it}\frac{\sin(\pi t)}{t}\widehat\kappa_{\rm nu}(t),
\]
where the quotient is interpreted continuously at zero. This proves
the endpoint zeros. Also
\begin{align*}
 \int_\pi^{3\pi}\frac{\sin s}{s}\,ds
 &=-\pi\int_0^\pi
       \frac{\sin u}{(\pi+u)(2\pi+u)}\,du
 \le-\frac1{3\pi}.
\end{align*}
Since $\sin s/s=\int_0^1\cos(ts)\,dt$ has derivative bounded by
$1/2$, convolution with $\kappa$ changes that integral by at most
$\pi\delta$. Therefore
\[
 K(0)=2\int_{\mathbb R}\frac{\sin s}{s}w(s)\,ds
 \le-\frac2{3\pi}+2\pi\delta<-\frac1{3\pi}.
\]
The last strict inequality follows from $6\pi^2<100$.
Twice integrating the definition of $K$ by parts, using $H(\pm1)=0$,
gives the claimed $u^{-2}$ decay; for example the numerator can be
bounded by $|H'(1)|+|H'(-1)|+\int_{-1}^1|H''|$.
Finally $H(-t)=\overline{H(t)}$, so $K(u)$ is real.
\end{proof}

\begin{proposition}[Conditional negative mass beside a critical zero]
\label{prop:v137-conditional-mass}
Assume RH. Fix any ordinate $\gamma_0$ of a nontrivial zero, of
multiplicity $m_0$, and use Lemma~\ref{lem:v137-probe}. Then
\begin{equation}\label{eq:v137-probe-limit}
 J_a:=\int_{\mathbb R}\beta_a(\xi)
              w\bigl(T(\xi-\gamma_0)\bigr)\,d\xi
 =m_0K(0)+O_{\gamma_0,w}(T^{-2})+E_T,
\end{equation}
where
\begin{equation}\label{eq:v137-exterior-error}
 |E_T|\le\frac{8\pi e^{-5T/2}}{5T(1-e^{-2T})}.
\end{equation}
Consequently, for all sufficiently large $a$,
\begin{equation}\label{eq:v137-conditional-lower}
 \Delta_a\ge\frac1{6\pi},\qquad
 \inf_{\xi\in\mathbb R}\beta_a(\xi)
      \le-\frac{a}{6\pi^2}.
\end{equation}
The threshold is not asserted effective here, and these displayed rates
are conditional on RH.
\end{proposition}
\begin{proof}
Write $\mathcal W$ for the complete geometric Weil distribution in
logarithmic coordinates. Under RH its spectral form is
$\mathcal W(t)=\sum_\gamma m_\gamma e^{i\gamma t}$ in the distributional
sense, with distinct real ordinates in this notation; this is the
classical explicit formula, not a new spectral assumption beyond RH
\cite[Theorem 2.1]{BurnolExplicit2000}.
The inverse-transform kernel of $\beta_a$ agrees with $\mathcal W$
on $|t|<T$. Outside that interval it is exactly
\[
 -\rho_*(|t|),\qquad
 \rho_*(t)=\frac{e^{-5t/2}}{1-e^{-2t}}\quad(t>0).
\]
Indeed this is the off-origin kernel of
$\Re\psi(5/4+i\xi/2)-\log\pi$. The prime and continuum terms in
\eqref{eq:v130-r-symbol} have support in $[-T,T]$; inside that interval
the identity
$-\rho_*(|t|)+e^{|t|/2}=-\rho(|t|)+2\cosh(t/2)$
retains the complete pole and archimedean terms. All distributions at
zero have the same normalization as \eqref{eq:v135-full-symbol}.

The nonunitary Fourier transform of the probe in \eqref{eq:v137-probe-limit}
is
\[
 V_T(t)=T^{-1}e^{-i\gamma_0t}H(t/T).
\]
The explicit formula applied to $\mathbf1_{[-T,T]}V_T$ therefore gives
the exact identity
\begin{equation}\label{eq:v137-zero-sum}
 J_a=\sum_\gamma m_\gamma K\bigl(T(\gamma-\gamma_0)\bigr)+E_T,
 \qquad
 E_T=-\int_{|t|>T}\rho_*(|t|)V_T(t)\,dt.
\end{equation}
There is no additional $2\pi$ in this identity: $V_T$ is the
nonunitary transform of the frequency test and $K$ is defined above
without a normalization factor. The ordinary physical Fourier transform
in \eqref{eq:v135-full-symbol} remains unitary.

For completeness the truncated test is legitimate. It is continuous,
compactly supported, smooth near zero and piecewise smooth elsewhere;
its endpoint values vanish because $H(\pm1)=0$. Its second distributional
derivative is a finite measure. Smooth mollification gives convergence
on the geometric side and Fourier decay $O((1+|\gamma|)^{-2})$
uniformly in the mollification parameter, for each fixed $T$.
The bound $N(Y)=O(Y\log(2+Y))$ then permits dominated convergence on
the RH zero sum. This extends the smooth explicit formula to this test.
If $e^T$ is a prime power, the endpoint atom pairs with zero; the strict
prime cutoff is therefore preserved exactly.

The ordinate $\gamma_0$ is isolated. Thus
\[
 \sum_{\gamma\ne\gamma_0}
       \frac{m_\gamma}{|\gamma-\gamma_0|^2}<\infty,
\]
and \eqref{eq:v137-probe-properties} bounds all other zero terms by
$O_{\gamma_0,w}(T^{-2})$. Since $|H(t)|\le2\pi$, integration of
$\rho_*$ bounds $E_T$ by \eqref{eq:v137-exterior-error}.
This proves \eqref{eq:v137-probe-limit}, and eventually
$J_a\le-1/(6\pi)$.

The scalar decomposition of Proposition~\ref{prop:v134-primitive}
gives $\beta_a=b_a+U_a'$ with $b_a\ge0$ and
$\norm{U_a}_\infty=\Delta_a/2$. Since the rescaled probe has variation
two, integration by parts yields
\[
 J_a\ge-\frac{\Delta_a}{2}
      \norm{(w(T(\cdot-\gamma_0)))'}_1=-\Delta_a.
\]
This proves the first inequality in \eqref{eq:v137-conditional-lower}.
Its mass is $2\pi/T=\pi/a$, which proves the second.
\end{proof}

\begin{corollary}[Unconditional scalar-budget obstruction]
\label{cor:v137-scalar-no-go}
For the exact arithmetic symbol of this manuscript,
\begin{equation}\label{eq:v137-scalar-divergence}
 a\Delta(\beta_a)\longrightarrow+\infty,
 \qquad
 \inf_{\xi\in\mathbb R}\beta_a(\xi)\longrightarrow-\infty
 \quad(a\longrightarrow\infty).
\end{equation}
In particular neither a bounded scalar primitive budget nor a bounded
pointwise negative error can supply the finite floor of
Proposition~\ref{prop:v136-bounded-floor}, even along a cofinal sequence.
No unconditional linear rate is claimed.
\end{corollary}
\begin{proof}
If $a_j\Delta(\beta_{a_j})$ were bounded on any sequence
$a_j\to\infty$, Proposition~\ref{prop:v134-primitive} would give a
uniform finite lower bound for the complete physical forms. The full
parity identity \eqref{eq:v135-full-symbol} and
Proposition~\ref{prop:v136-bounded-floor} would then imply RH.
But under RH, Proposition~\ref{prop:v137-conditional-mass} forces
$a_j\Delta(\beta_{a_j})\ge a_j/(6\pi)$ eventually, a contradiction.
Absence of a bounded cofinal subsequence is precisely the first limit
in \eqref{eq:v137-scalar-divergence}.
Likewise a cofinal pointwise lower bound $\beta_{a_j}\ge-C$ would
give a complete physical lower floor, imply RH, and contradict the
second estimate in \eqref{eq:v137-conditional-lower}.
This proves the second limit.
\end{proof}

This is not a negative physical Weil test. The nonnegative frequency
probe $w(T(\xi-\gamma_0))$ is not asserted to be the squared transform
of a vector supported in $(-a,a)$. Such an assertion would contradict
the RH positivity used in the conditional part of the argument.
The result identifies exactly the cost of discarding physical
time--frequency correlations. Both parities remain within the complete
symbol identity, but no projection, low-concentration identification,
source-to-physical transfer or sampler graph formula is changed.

The nearest project proposal is the v1.34 scalar primitive estimate,
retested under v1.36's weaker bounded-error objective. The new ingredient
is an endpoint-vanishing positive frequency probe, which gives an
absolutely convergent zero sum and a surviving negative cutoff side lobe.
This is a proved obstruction to that route, not a new positive mechanism.
The explicit formula and cutoff oscillations are classical. Nearby
pointwise-envelope barriers, for example \cite{Zhu2026}, bound the raw
prime comb; the present statement concerns the continuum-cancelled
symbol and its optimal signed primitive. No worldwide novelty is claimed.
The joint signed concentration estimate remains open: it must retain
the favorable levels jointly with their physical concentration operators
to obtain a uniform finite ordinary lower error. No G2 sign gap is closed.

